% !TeX root = ../main.tex
% =========================================================================
% 8. CRONOGRAMA DE ACTIVIDADES (400 horas + 10 % adicional si se requiere)
% Formato de tablas APA 7 estricto (booktabs, sin líneas verticales)
% =========================================================================
\section{Cronograma de actividades}

El desarrollo del proyecto de titulación se estructurará a lo largo de un período académico regular, ajustándose de forma estricta a los hitos de evaluación y seguimiento dispuestos por la Comisión de Titulación de la carrera, tal como se detalla en la Tabla \ref{tab:cronograma}.

\begin{table}[H]
\centering
\caption{Cronograma de hitos del programa de titulación}
\label{tab:cronograma}
\begin{tabular}{p{5.0cm} p{7.0cm} c}
\toprule
\textbf{Hito institucional} & \textbf{Entregable / Alcance} & \textbf{Fecha límite} \\
\midrule
Presentación y aprobación & Aprobación oficial del tema de titulación & 30/11/2026 \\
Hito 1 (25\%) & Estado del arte, requerimientos y catálogos & 21/12/2026 \\
Hito 2 (50\%) & Diseño arquitectónico y esquema de datos & 25/01/2027 \\
Hito 3 (75\%) & Prototipo funcional desplegado en la nube & 01/03/2027 \\
Hito 4 (100\%) & Evaluación experimental e informe final & 05/04/2027 \\
\bottomrule
\end{tabular}
\notatabla{Cronograma oficial de hitos académicos y ponderaciones de avance de acuerdo con la planificación de la Universidad Politécnica Salesiana.}
\end{table}

Para garantizar la culminación exitosa de los objetivos planteados dentro del límite normativo de 400 horas de dedicación estudiantil, se presenta en la Tabla \ref{tab:actividades} el desglose pormenorizado de actividades por cada una de las cinco fases de la investigación, indicando su carga horaria estimada, fechas de inicio y culminación, y los responsables asignados.

\begin{table}[H]
\centering
\caption{Plan de actividades por fases del proyecto de titulación}
\label{tab:actividades}
\small
\begin{tabular}{p{5.2cm} c c c c}
\toprule
\textbf{Actividad} & \textbf{Horas} & \textbf{Inicio} & \textbf{Fin} & \textbf{Responsables} \\
\midrule
\multicolumn{5}{l}{\textit{Fase 1: Revisión bibliográfica y análisis de requerimientos}} \\
Revisión sistemática de literatura científica & 40 & 01/12/2026 & 12/12/2026 & Autores \\
Levantamiento de requerimientos y digitalización & 30 & 14/12/2026 & 21/12/2026 & Autores / Tutor \\
\addlinespace
\multicolumn{5}{l}{\textit{Fase 2: Diseño de la arquitectura propuesta}} \\
Diseño arquitectónico y modelado relacional SQLite & 40 & 22/12/2026 & 11/01/2027 & Autores \\
Estructuración de índices y normalización multimarca & 30 & 12/01/2027 & 25/01/2027 & Autores \\
\addlinespace
\multicolumn{5}{l}{\textit{Fase 3: Desarrollo del prototipo funcional}} \\
Implementación del backend .NET 8 y memoria de sesión & 60 & 26/01/2027 & 12/02/2027 & Autores \\
Integración con API Gemini y WhatsApp Cloud API & 40 & 15/02/2027 & 01/03/2027 & Autores \\
\addlinespace
\multicolumn{5}{l}{\textit{Fase 4: Despliegue y pruebas experimentales}} \\
Despliegue en AWS con Nginx TLS y batería de pruebas & 60 & 02/03/2027 & 19/03/2027 & Autores \\
\addlinespace
\multicolumn{5}{l}{\textit{Fase 5: Análisis de resultados y documentación}} \\
Análisis estadístico de telemetría y redacción del informe & 60 & 20/03/2027 & 05/04/2027 & Autores / Tutor \\
\midrule
Subtotal de horas planificadas & 360 & & & \\
Horas de contingencia y ajustes (10\%) & 40 & & & \\
\midrule
\textbf{Total planificado} & \textbf{400} & & & \\
\bottomrule
\end{tabular}
\notatabla{Distribución del esfuerzo de ingeniería en 360 horas de ejecución técnica y 40 horas de contingencia para imprevistos o ajustes experimentales, completando las 400 horas reglamentarias.}
\end{table}

La articulación entre el plan de actividades (Tabla \ref{tab:actividades}) y los hitos institucionales (Tabla \ref{tab:cronograma}) asegura un avance disciplinado y continuo, permitiendo la verificación temprana de los componentes de software antes de la fase de evaluación experimental definitiva.
